\documentclass{article}

\usepackage{fancyhdr}
\usepackage{extramarks}
\usepackage{amsmath}
\usepackage{amsthm}
\usepackage{amsfonts}
\usepackage{tikz}
\usepackage[plain]{algorithm}
\usepackage{algpseudocode}
\usepackage{fullpage, url, hyperref}
\newcommand{\dotcup}{\ensuremath{\mathaccent\cdot\cup}}

\theoremstyle{plain}
\newtheorem{theorem}{Theorem}
\newtheorem{lemma}{Lemma}
\newtheorem{definition}{Definition}

\newcommand{\I}{\mathcal{I}}
\newcommand{\F}{\mathbb{F}}
\newcommand{\E}{\mathbb{E}}
\newcommand{\R}{\mathbb{R}}
\newcommand{\tr}{\mathrm{tr}}
\renewcommand{\P}{\mathbb{P}}
\usetikzlibrary{automata,positioning}

%
% Basic Document Settings
%

\topmargin=-0.45in
\evensidemargin=-.50in
\oddsidemargin=-.5in
\textwidth=7in
\textheight=9.5in
\headsep=0.25in

\linespread{1.1}

\pagestyle{fancy}
\lhead{\hmwkAuthorName}
\chead{\hmwkClass\ (\hmwkClassInstructor\ \hmwkClassTime): \hmwkTitle}
\rhead{\firstxmark}
\lfoot{\lastxmark}
\cfoot{\thepage}

\newcommand{\w}[2]{w_{#1}^{(#2)}}
\newcommand{\loss}[2]{\vec{\ell}_{#2}^{(#1)}}
\newcommand{\losst}[1]{\vec{\ell}^{(#1)}}
  \newcommand{\pt}[1]{\vec{p}^{(#1)}}
\renewcommand\headrulewidth{0.4pt}
\renewcommand\footrulewidth{0.4pt}

\setlength\parindent{0pt}


% 
% Create Problem Sections
%


\newcommand{\enterProblemHeader}[1]{
    \nobreak\extramarks{}{}\nobreak{}
    \nobreak\extramarks{}\nobreak{}
}

\newcommand{\exitProblemHeader}[1]{
    \nobreak\extramarks{}\nobreak{}
    \stepcounter{#1}
    \nobreak\extramarks{}\nobreak{}
  }

  
\newcommand{\enterExerciseHeader}[1]{
    \nobreak\extramarks{}{}\nobreak{}
    \nobreak\extramarks{}\nobreak{}
}

\newcommand{\exitExerciseHeader}[1]{
    \nobreak\extramarks{}\nobreak{}
    \stepcounter{#1}
    \nobreak\extramarks{}\nobreak{}
}

\setcounter{secnumdepth}{0}
\newcounter{partCounter}
\newcounter{homeworkProblemCounter}
\setcounter{homeworkProblemCounter}{1}

\newcounter{homeworkExerciseCounter}
\setcounter{homeworkExerciseCounter}{1}
\nobreak\extramarks{}\nobreak{}
\nobreak\extramarks{}\nobreak{}

%
% Homework Problem Environment
%
% This environment takes an optional argument. When given, it will adjust the
% problem counter. This is useful for when the problems given for your
% assignment aren't sequential. See the last 3 problems of this template for an
% example.
%
\newenvironment{homeworkProblem}[1][-1]{
%    \ifnum#1>0
   %     \setcounter{homeworkProblemCounter}{#1}
   % \fi
    \section{Problem \arabic{homeworkProblemCounter}}
    \setcounter{partCounter}{1}
    \enterProblemHeader{homeworkProblemCounter}
}{
    \exitProblemHeader{homeworkProblemCounter}
}

  \newenvironment{homeworkExercise}[1][-1]{
    \ifnum#1>0
        \setcounter{homeworkExerciseCounter}{#1}
    \fi
    \section{Exercise \arabic{homeworkExerciseCounter}}
    \setcounter{partCounter}{1}
    \enterExerciseHeader{homeworkExerciseCounter}
}{
    \exitExerciseHeader{homeworkExerciseCounter}
}

%
% Homework Details
%   - Title
%   - Due date
%   - Class
%   - Section/Time
%   - Instructor
%   - Author
%

\newcommand{\hmwkTitle}{Homework\ \#1}
\newcommand{\hmwkAssignedDate}{Aug 28}
\newcommand{\hmwkDueDate}{Sept 10, midnight}
\newcommand{\hmwkClass}{Foundations of Fairness in Machine Learning, Fall 2018}
\newcommand{\hmwkClassTime}{}%Section A}
\newcommand{\hmwkClassInstructor}{Professor Jamie Morgenstern}
\newcommand{\hmwkAuthorName}{}%\textbf{Josh Davis} \and \textbf{Davis Josh}}

%
% Title Page
%

\title{    \vspace{2in}
    \textmd{\textbf{\hmwkClass:\ \hmwkTitle}}\\
    \normalsize\vspace{0.1in}\small{Due\ on\ \hmwkDueDate\ at 3:10pm}\\
    \vspace{0.1in}\large{\textit{\hmwkClassInstructor\ \hmwkClassTime}}
    \vspace{3in}
}

\author{\hmwkAuthorName}
\date{}

\renewcommand{\part}[1]{\textbf{\large Part \Alph{partCounter}}\stepcounter{partCounter}\\}

%
% Various Helper Commands
%

% Useful for algorithms
\newcommand{\alg}[1]{\textsc{\bfseries \footnotesize #1}}

% For derivatives
\newcommand{\deriv}[1]{\frac{\mathrm{d}}{\mathrm{d}x} (#1)}

% For partial derivatives
\newcommand{\pderiv}[2]{\frac{\partial}{\partial #1} (#2)}

% Integral dx
\newcommand{\dx}{\mathrm{d}x}

% Alias for the Solution section header
\newcommand{\solution}{\textbf{\large Solution}}

% Probability commands: Expectation, Variance, Covariance, Bias
\newcommand{\Var}{\mathrm{Var}}
\newcommand{\Cov}{\mathrm{Cov}}
\newcommand{\Bias}{\mathrm{Bias}}
\usepackage{enumitem}
\newcommand\litem[1]{\item{\bfseries#1.\space}}

\newcommand{\Z}{\mathbb{Z}}

\newcommand{\poly}{\textrm{poly}}
\begin{document}

% \maketitle

\pagebreak

{\bf Homework Out: \hmwkAssignedDate}\\
{\bf Due Date:  \hmwkDueDate}\\
{\bf Reminder: }

To submit your homework, please go to
\url{https://classroom.github.com/a/7rslWmMd}, accept the assignment,
and submit your LaTex, PDF, and any code you use for the
assignment. Please name your files ``hw1-USERNAME-writeup.{tex,pdf}''
``hw1-USERNAME-code.{appropriate file type}''. You will need a github
account, and to add, commit, and push your homework.

Please cite all sources you use, and people you work with. The
expectation is that you try and solve these problems yourself, rather
than looking online explicitly for answers. Submissions due at 23:00
of the due date.

{\bf You may use $O$-notation unless explicitly noted somewhere in the homework.}


\newcommand{\X}{\mathcal{X}}
\newcommand{\Y}{\mathcal{Y}}
\newcommand{\D}{\mathcal{D}}

\paragraph{Problems}
\begin{enumerate}
\litem{(Statistical Guarantees aren't Created Equal.)}


Consider some $f : \X \to [0,1]$ and distributions $\D_1, \D_2$ over
$\X$ such that

\[\E_{x \sim \D_1} [f(x)] = \E_{x \sim \D_2} [f(x)]\]

Consider two datasets $X_1, X_2$ drawn independently as
$ x \sim \D_i, x \in X_i$, with $|X_1| = n_1 > n_2 = |X_2|$.

\begin{enumerate}
  \item
As a
function of $n_1, n_2$, give an upper bound on

\[ |\frac{1}{n_1} \sum_{x \in X_1} f(x) - \frac{1}{n_2}\sum_{x \in X_2} f(x)|\]
%
which holds with probability $1-\delta$ over the draw of $X_1,
X_2$. You can use ``known'' concentration bounds (e.g., Chernoff-Hoeffding).

\item Lower bound the probability with which
  \[\frac{1}{n_1} \sum_{x \in X_1} f(x)   \geq  \frac{1}{n_2}\sum_{x \in X_2} f(x) + \frac{1}{\sqrt{n_2}}. \]

  Notice this describes the probability that two empirical averages differ by the \emph{larger} of their variances.
\end{enumerate}


\litem{(Distortion vs Information Theory.)}


The paper we read (``On the (Im)possibility of Fairness'') for the
first class suggested we consider distortion as a measure of WYSIWYG
(how much our observed features represent the construct space for
different groups). Let $\X, d_{\X}$ represent the construct space and
its metric, and $\Y, d_{\Y}$ the observed space and metric,
respectively.

\begin{enumerate}
  \item Describe some $\X\subseteq \R^n$ and $f$ comprised of some linear
  measurements of $\X$, $f$ transforming $\X$ to $\Y \subseteq \R^{n'}$ such that the
  distortion of $f$ is of order $\max_{x \in \X}||x ||_2$, but the
  accuracy of the \emph{best} linear classifier is no worse in $\Y$
  than in $\X$.

  You can assume $d_x = d_y$ is the $\ell_2$ metric. 

  \item Find some $\X$, $f$, and two groups $\X_1\dotcup \X_2 = \X$
  such that the distortion of $f$ on $\X_1$ (ignoring points in $\X_2$) and
  the distortion of $\X_2$ (ignoring points in $\X_1$) are equal, but
  the best linear classifier for $\X$ has $0$ error, the best linear
  classifier for $f(\X_1)$ has error $\frac{1}{2}$, and the best
  linear classifier for $f(\X_2)$ has error $0$.
  \end{enumerate}



\litem{(Coding: Training
  together, training apart.)}

Find a human-centric dataset\footnote{A dataset such that each entry
  represents one measurement of one person.} which is publicly
available with demographic information (age, gender, race, sexual
orientation, country of origin, etc).

\begin{enumerate}
\item 

Please include a link to the dataset you used, as well as any
documentation that accompanied its release. Clearly describe any
``cleaning'', binning, bucketing, or discrete-to-continuous feature
transformation you did in this process.

\item For one way of splitting the dataset into different demographic
  groups, write down the size of the groups, the average value for
  each (numeric) feature, the variance of each (numeric) feature, the
  mode for each categorical feature, and the three most frequent
  values for each categorical feature, each computed on the different demographic subgroups.

  If your dataset has more than 20 features, you can report this information only for 20 features.

  Do you observe any interesting differences between the different subgroups' statistics?

\item Now, randomly subsample half of the dataset. That is, include
  each row in the database in some new dataset with probability
  $\frac{1}{2}$ independently.  Report the same statistics as above.

  Do you observe some of the statistics are more ``stable'' than others (namely, that their values
  changed little from the previous question)? Do some subgroups have more stable statistics than others?

  

  \end{enumerate}



  \end{enumerate}

\end{document}

%%% Local Variables:
%%% mode: latex
%%% TeX-master: t
%%% End:
